\documentclass{article}
\usepackage[margin=1in]{geometry}
\usepackage{hyperref}
\usepackage{ulem}
\usepackage{tipa}
\hypersetup{
    colorlinks=true,
    linkcolor=black,
    filecolor=magenta,
    urlcolor=cyan,
}

\urlstyle{same}

\title{French Grammar and Connectors Summary}
\author{Junzhang Luo}

\begin{document}

\tableofcontents

\newpage

\section{Present Participle and Gerund}

\subsection{Participe présent}
\begin{itemize}
    \item \textbf{Core use:} Adds descriptive information (similar to a relative clause) in formal written French.
    \item \textbf{Formation / pattern:} Take the \textit{nous} form in present tense, remove \textit{-ons}, add \textit{-ant} (\textit{nous parlons} $\to$ \textit{parlant}).
    \item \textbf{Key contrast:} \textit{Participe présent} describes a noun; it is not the same as the English progressive tense.
    \item \textbf{Example:} \textit{Les étudiants parlant français trouvent plus facilement un stage.} (English: Students who speak French find an internship more easily.)
\end{itemize}

\subsection{Gérondif}
\begin{itemize}
    \item \textbf{Core use:} Expresses simultaneous action, means, or cause.
    \item \textbf{Formation / pattern:} \textit{en + participe présent}.
    \item \textbf{Key contrast:} \textit{Gérondif} usually describes the subject's own action in context (\textit{en travaillant}, \textit{en cherchant}).
    \item \textbf{Example:} \textit{Il a appris le français en regardant des films.} (English: He learned French by watching films.)
\end{itemize}

\section{Adverb Formation}
\begin{itemize}
    \item \textbf{Core use:} Adverbs modify verbs, adjectives, or entire statements.
    \item \textbf{Formation / pattern:} Usually feminine adjective form + \textit{-ment} (\textit{rapide} $\to$ \textit{rapidement}).
    \item \textbf{Key contrast:} \textit{-ant} adjectives often become \textit{-amment}; \textit{-ent} adjectives often become \textit{-emment}.
    \item \textbf{Example:} \textit{Elle parle couramment français.} (English: She speaks French fluently.)
\end{itemize}

\section{Imparfait}
\begin{itemize}
    \item \textbf{Core use:} Background, habits, repeated actions, and states in the past.
    \item \textbf{Formation / pattern:} Stem from \textit{nous} present minus \textit{-ons} + endings \textit{-ais, -ais, -ait, -ions, -iez, -aient}.
    \item \textbf{Key contrast:} \textit{être} is irregular in stem: \textit{ét-}.
    \item \textbf{Example:} \textit{Quand j'étais petit, je lisais tous les soirs.} (English: When I was little, I read every evening.)
\end{itemize}

\section{Plus-que-parfait}
\begin{itemize}
    \item \textbf{Core use:} Action completed before another past action.
    \item \textbf{Formation / pattern:} Imperfect of auxiliary (\textit{avoir} or \textit{être}) + past participle.
    \item \textbf{Key contrast:} Use it to establish a "past before past" timeline with \textit{passé composé} or \textit{imparfait}.
    \item \textbf{Example:} \textit{Il avait déjà dîné quand je suis arrivé.} (English: He had already had dinner when I arrived.)
\end{itemize}

\section{Adjectives}

\subsection{Possessive adjectives}
\begin{itemize}
    \item \textbf{Core use:} Shows ownership or relation.
    \item \textbf{Formation / pattern:} \textit{mon/ma/mes, ton/ta/tes, son/sa/ses, notre/nos, votre/vos, leur/leurs}.
    \item \textbf{Key contrast:} Choose according to the gender/number of the noun possessed, not the owner's gender.
    \item \textbf{Example:} \textit{Marie cherche son ordinateur.} (English: Marie is looking for her computer.)
\end{itemize}

\subsection{Superlatives}
\begin{itemize}
    \item \textbf{Core use:} Expresses highest or lowest degree.
    \item \textbf{Formation / pattern:} \textit{le/la/les plus + adjective} or \textit{le/la/les moins + adjective}.
    \item \textbf{Key contrast:} Irregular forms include \textit{le meilleur / la meilleure} and \textit{le pire / la pire}.
    \item \textbf{Example:} \textit{C'est la meilleure solution pour ce projet.} (English: It's the best solution for this project.)
\end{itemize}

\subsection{Adjective gender-change patterns}
\begin{itemize}
    \item \textbf{Core use:} Build feminine adjective forms from masculine base forms.
    \item \textbf{Formation / pattern:} See common endings below.
    \item \textbf{Key contrast:} Some adjectives are irregular (\textit{beau $\to$ belle}, \textit{nouveau $\to$ nouvelle}, \textit{vieux $\to$ vieille}); some are invariable in spoken usage.
    \item \textbf{Example:} \textit{Ce quartier est ancien, mais cette rue est ancienne.} (English: This district is old, but this street is old.)
\end{itemize}
\begin{center}
    \begin{tabular}{|l|l|l|}
        \hline
        \textbf{Pattern} & \textbf{Masculine $\to$ Feminine} & \textbf{Example pair} \\
        \hline
        -el $\to$ -elle & cruel $\to$ cruelle & un ton cruel / une remarque cruelle \\
        \hline
        -if $\to$ -ive & sportif $\to$ sportive & un élève sportif / une élève sportive \\
        \hline
        -eux $\to$ -euse & heureux $\to$ heureuse & un enfant heureux / une enfant heureuse \\
        \hline
        -er $\to$ -ère & cher $\to$ chère & un livre cher / une voiture chère \\
        \hline
        -f $\to$ -ve & neuf $\to$ neuve & un sac neuf / une robe neuve \\
        \hline
        -c $\to$ -que & public $\to$ publique & un débat public / une décision publique \\
        \hline
        -on $\to$ -onne & bon $\to$ bonne & un bon repas / une bonne idée \\
        \hline
        -en $\to$ -enne & ancien $\to$ ancienne & un ancien quartier / une ancienne rue \\
        \hline
        -et $\to$ -ette / -ète & net $\to$ nette; complet $\to$ complète & un texte net / une réponse nette \\
        \hline
    \end{tabular}
\end{center}

\section{Future}

\subsection{Futur simple}
\begin{itemize}
    \item \textbf{Core use:} Future facts, plans, predictions, and formal instructions.
    \item \textbf{Formation / pattern:} Infinitive (or irregular stem) + endings \textit{-ai, -as, -a, -ons, -ez, -ont}.
    \item \textbf{Key contrast:} Common irregular stems: \textit{ser-, aur-, ir-, fer-, viendr-, verr-, pourr-, voudr-, devr-}.
    \item \textbf{Example:} \textit{Nous verrons les résultats demain.} (English: We will see the results tomorrow.)
\end{itemize}

\section{Pronouns}

\subsection{Direct object pronouns}
\begin{itemize}
    \item \textbf{Core use:} Replaces direct objects.
    \item \textbf{Formation / pattern:} \textit{me, te, le, la, nous, vous, les}.
    \item \textbf{Key contrast:} Usually before the conjugated verb (except affirmative imperative).
    \item \textbf{Example:} \textit{Je le connais bien.} (English: I know him well.)
\end{itemize}

\subsection{Indirect object pronouns}
\begin{itemize}
    \item \textbf{Core use:} Replaces objects introduced by \textit{à} (to someone).
    \item \textbf{Formation / pattern:} \textit{me, te, lui, nous, vous, leur}.
    \item \textbf{Key contrast:} Use indirect pronouns with verbs that take \textit{à} + person.
    \item \textbf{Example:} \textit{Je lui téléphone ce soir.} (English: I'm calling him/her tonight.)
\end{itemize}

\subsection{Relative pronouns}
\begin{itemize}
    \item \textbf{Core use:} Relie deux propositions et évite la répétition.
    \item \textbf{Formation / pattern:} \textit{qui} (sujet), \textit{que} (COD), \textit{dont} (complément avec \textit{de}), \textit{où} (lieu/temps).
    \item \textbf{Key contrast:} Pour les choix détaillés (\textit{qui/que/dont/où}, \textit{lequel}, \textit{auquel}, \textit{duquel}), voir la section ``When to use''.
    \item \textbf{Example:} \textit{La ville où j'habite est celle dont je t'ai parlé.} (English: The city where I live is the one I told you about.)
\end{itemize}

\subsection{\textit{y} and \textit{en}}
\begin{itemize}
    \item \textbf{Core use:} Compact replacements for frequent prepositional complements.
    \item \textbf{Formation / pattern:} \textit{y} replaces \textit{à} + thing/place; \textit{en} replaces \textit{de} + thing/quantity.
    \item \textbf{Key contrast:} With quantity, keep the number after \textit{en} (\textit{J'en veux deux}).
    \item \textbf{Example:} \textit{Tu vas au marché? Oui, j'y vais.} (English: Are you going to the market? Yes, I'm going.)
\end{itemize}

\subsection{Stressed pronouns}
\begin{itemize}
    \item \textbf{Core use:} Emphasis, contrast, and use after prepositions.
    \item \textbf{Formation / pattern:} \textit{moi, toi, lui, elle, nous, vous, eux, elles}.
    \item \textbf{Key contrast:} Often used in oral emphasis: \textit{Moi, je pense que...}
    \item \textbf{Example:} \textit{Ce cadeau est pour elle.} (English: This gift is for her.)
\end{itemize}

\section{Negation}
\begin{itemize}
    \item \textbf{Core use:} Denies, limits, or restricts information.
    \item \textbf{Formation / pattern:} Basic frame \textit{ne ... pas}; common variants: \textit{ne ... jamais}, \textit{ne ... rien}, \textit{ne ... personne}, \textit{ne ... plus}, \textit{ne ... que}.
    \item \textbf{Key contrast:} In spoken French, \textit{ne} is often dropped, but keep full form in writing.
    \item \textbf{Example:} \textit{Personne ne comprend ce message.} (English: No one understands this message.)
\end{itemize}

\section{Conditionnel}

\subsection{Conditionnel présent}
\begin{itemize}
    \item \textbf{Core use:} Polite requests, wishes, advice, hypotheses, and soft assertions.
    \item \textbf{Formation / pattern:} Futur stem + imparfait endings \textit{-ais, -ais, -ait, -ions, -iez, -aient}.
    \item \textbf{Key contrast:} In hypotheses, use \textit{si + imparfait} $\to$ \textit{conditionnel présent}.
    \item \textbf{Example:} \textit{Si j'avais plus de temps, je voyagerais davantage.} (English: If I had more time, I would travel more.)
\end{itemize}

\subsection{Conditionnel passé}
\begin{itemize}
    \item \textbf{Core use:} Past hypothesis, regret, reproach, and unconfirmed information (often journalistic style).
    \item \textbf{Formation / pattern:} Auxiliary in conditionnel présent (\textit{avoir} or \textit{être}) + past participle.
    \item \textbf{Key contrast:} With \textit{si}, use \textit{si + plus-que-parfait} $\to$ \textit{conditionnel passé}; with \textit{être}, participle agrees in gender/number.
    \item \textbf{High-frequency participles:} \textit{pouvoir} $\to$ \textit{pu}, \textit{devoir} $\to$ \textit{dû}, \textit{venir} $\to$ \textit{venu(e)(s)}, \textit{vouloir} $\to$ \textit{voulu}, \textit{savoir} $\to$ \textit{su}.
    \item \textbf{Example:} \textit{J'aurais dû vérifier les horaires avant de partir.} (English: I should have checked the schedules before leaving.)
\end{itemize}

\section{Subjonctif}
\begin{itemize}
    \item \textbf{Core use:} Emotion, doubt, necessity, judgment, and subjectivity after \textit{que}.
    \item \textbf{Formation / pattern:} Stem + endings \textit{-e, -es, -e, -ions, -iez, -ent}; many \textit{-ir} verbs insert \textit{-iss-}.
    \item \textbf{Key contrast:} Frequent irregulars: \textit{être} (\textit{sois, sois, soit, soyons, soyez, soient}) and \textit{avoir} (\textit{aie, aies, ait, ayons, ayez, aient}).
    \item \textbf{Example:} \textit{Il faut que nous soyons prêts avant midi.} (English: We must be ready before noon.)
\end{itemize}

\section{Impératif}
\begin{itemize}
    \item \textbf{Core use:} Commands, instructions, invitations, and advice.
    \item \textbf{Formation / pattern:} Uses \textit{tu, nous, vous} forms.
    \item \textbf{Key contrast:} In \textit{-er} verbs, final \textit{-s} usually drops in affirmative \textit{tu} form (except before \textit{y/en}); irregulars include \textit{sois/soyons/soyez}, \textit{aie/ayons/ayez}, \textit{sache/sachons/sachez}, \textit{veuille/veuillons/veuillez}.
    \item \textbf{Example:} \textit{Prenons le temps de bien relire le texte.} (English: Let's take time to carefully reread the text.)
\end{itemize}

\section{Reported Speech}
\begin{itemize}
    \item \textbf{Core use:} Reports statements/questions from another speaker.
    \item \textbf{Formation / pattern:} Common reporting verbs: \textit{dire que}, \textit{affirmer que}, \textit{demander si}.
    \item \textbf{Key contrast:} With tense shift in a past reporting frame: present $\to$ imparfait, futur $\to$ conditionnel, \textit{aller + infinitif} (near future) $\to$ imparfait of \textit{aller} + infinitive.
    \item \textbf{Example:} \textit{Il a dit qu'il viendrait le lendemain.} (English: He said that he would come the next day.)
\end{itemize}

\section{Nominalisation}
\begin{itemize}
    \item \textbf{Core use:} Converts actions/qualities into nouns for denser and more formal style.
    \item \textbf{Formation / pattern:} Verb $\to$ noun (\textit{-ment, -tion, -ture, -ée}); adjective $\to$ noun (\textit{-tude, -eur, -ance/-ence, -esse}).
    \item \textbf{Key contrast:} Useful in formal writing, but overuse can reduce clarity in spoken style.
    \item \textbf{Example:} \textit{La modernisation du système améliore la communication.} (English: Modernizing the system improves communication.)
\end{itemize}

\section{Connecteurs et vocabulaire utile}

\subsection{But / objectif}
\begin{itemize}
    \item \textit{afin de} (in order to)
    \item \textit{pour} (to, in order to)
    \item \textit{avoir pour fonction de + infinitif} (to have the role/function of)
\end{itemize}
\textbf{Exemple modèle:} \textit{Afin d'améliorer son niveau, elle lit chaque jour pour progresser plus vite.} (English: To improve her level, she reads every day to make faster progress.)

\subsection{Cause}
\begin{itemize}
    \item \textit{car} (because)
    \item \textit{parce que} (because)
    \item \textit{grâce à} (thanks to)
    \item \textit{en raison de} (due to)
    \item \textit{comme} / \textit{et comme} (as, since)
\end{itemize}
\textbf{Exemple modèle:} \textit{Comme il pleuvait, nous sommes restés à la maison en raison du mauvais temps.} (English: Since it was raining, we stayed at home because of the bad weather.)

\subsection{Conséquence}
\begin{itemize}
    \item \textit{alors} (so, then)
    \item \textit{donc} (therefore, so)
    \item \textit{du coup} (so, as a result)
    \item \textit{ainsi} (thus)
    \item \textit{par conséquent} (consequently)
\end{itemize}
\textbf{Exemple modèle:} \textit{Le train a été annulé, donc nous avons pris un taxi; par conséquent, nous sommes arrivés en retard.} (English: The train was canceled, so we took a taxi; consequently, we arrived late.)

\subsection{Opposition / concession}
\begin{itemize}
    \item \textit{malgré} (despite)
    \item \textit{même si} (even if)
    \item \textit{en revanche} (on the other hand)
    \item \textit{par contre} (on the other hand)
    \item \textit{néanmoins} (nevertheless)
    \item \textit{toutefois} (however)
    \item \textit{cependant} (however)
    \item \textit{il n'empêche que} (the fact remains that)
    \item \textit{contre} (against)
    \item \textit{tout en} (while)
\end{itemize}
\textbf{Exemple modèle:} \textit{Même si le plan est ambitieux, il n'empêche que le budget reste limité.} (English: Even if the plan is ambitious, the fact remains that the budget is limited.)

\subsection{Addition / renforcement}
\begin{itemize}
    \item \textit{d'ailleurs} (besides, moreover)
    \item \textit{en outre} (furthermore)
    \item \textit{en plus} (in addition)
    \item \textit{par ailleurs} (moreover)
    \item \textit{surtout} (especially)
    \item \textit{tout à fait} (absolutely)
    \item \textit{Il n'y a pas que ... mais aussi ...} (not only ... but also ...)
\end{itemize}
\textbf{Exemple modèle:} \textit{Ce cours est utile et, en plus, il n'y a pas que la grammaire mais aussi la prononciation.} (English: This course is useful and, in addition, it covers not only grammar but also pronunciation.)

\subsection{Illustration / exemple}
\begin{itemize}
    \item \textit{par exemple} (for example)
    \item \textit{un de ces + nom} (an impressive/emphatic noun phrase)
    \item \textit{personnages hauts en couleurs} (colorful characters)
\end{itemize}
\textbf{Exemple modèle:} \textit{Par exemple, ce roman présente un de ces héros hauts en couleurs qu'on n'oublie pas.} (English: For example, this novel features one of those vivid heroes you don't forget.)

\subsection{Reformulation / précision}
\begin{itemize}
    \item \textit{c'est-à-dire} (that is to say)
    \item \textit{à savoir} (namely)
    \item \textit{en fait} (in fact)
    \item \textit{en effet} (indeed)
    \item \textit{enfin} (I mean, rather)
    \item \textit{enfin surtout} (well, mostly)
    \item \textit{plutôt} (rather)
    \item \textit{vouloir dire} (to mean)
    \item \textit{qu'en est-il de ... ?} (what about ...?)
\end{itemize}
\textbf{Exemple modèle:} \textit{Le projet avance, c'est-à-dire que les objectifs principaux sont déjà atteints.} (English: The project is moving forward, that is to say, the main objectives have already been met.)

\subsection{Résumé / conclusion}
\begin{itemize}
    \item \textit{en gros} (basically)
    \item \textit{en résumé} (in summary)
    \item \textit{en conclusion} (in conclusion)
    \item \textit{ce qui compte} (what matters)
\end{itemize}
\textbf{Exemple modèle:} \textit{En résumé, ce qui compte est de pratiquer régulièrement.} (English: In summary, what matters is practicing regularly.)

\subsection{Point de vue / cadrage}
\begin{itemize}
    \item \textit{de mon côté} (as for me)
    \item \textit{selon} (according to)
    \item \textit{à l'égard de} (with regard to)
    \item \textit{à propos de} (about)
    \item \textit{en matière de} (in terms of)
    \item \textit{ayant trait à} (relating to)
    \item \textit{en provenance de} (coming from)
    \item \textit{au-delà de} (beyond)
    \item \textit{au lieu de} (instead of)
    \item \textit{aller dans un sens} (to point in a direction)
\end{itemize}
\textbf{Exemple modèle:} \textit{De mon côté, en matière de lecture, je choisis des textes ayant trait à la culture francophone.} (English: For my part, in terms of reading, I choose texts related to Francophone culture.)

\subsection{Quantité / comparaison}
\begin{itemize}
    \item \textit{la plupart de} (most of)
    \item \textit{autant de} (as many)
    \item \textit{autant que} (as much/as far as)
    \item \textit{tant de} (so much, so many)
    \item \textit{parmi} (among)
    \item \textit{être + nombre/pourcentage + à + infinitif} (to be at a given number/percentage in doing)
\end{itemize}
\textbf{Exemple modèle:} \textit{La plupart des étudiants lisent autant que possible pour progresser.} (English: Most students read as much as possible to improve.)

\subsection{Expressions figées utiles}
\begin{itemize}
    \item \textit{y aller de} (to contribute/to come out with)
    \item \textit{Vivement que + subjonctif} (I can't wait for ...)
    \item \textit{ne + verbe + rien} (nothing)
    \item \textit{ne ... plus} (no longer)
    \item \textit{ne ... plus que} (only remains)
    \item \textit{ne ... donc plus que} (thus only remains)
    \item \textit{répéter \textit{que} pour deux propositions:} \textit{Il a dit que ..., et que ...}
    \item \textit{permettre de + infinitif} (to allow)
    \item \textit{les uns et les autres} (one another)
    \item \textit{quoi que ce soit} (anything at all)
    \item \textit{venir de + infinitif} (to have just done)
    \item \textit{Il arrive que ...} (it happens that ...)
\end{itemize}
\textbf{Exemple modèle:} \textit{Vivement qu'il arrive, car je ne peux plus attendre.} (English: I can't wait for him to arrive, because I can't wait any longer.)

\section{When to use: pronoms relatifs et prépositions fréquentes}

\subsection{De vs Des}
    \begin{itemize}
        \item When talking in general, use de.
        \item When talking specifc (of the ...) use all the forms of de dependent on the noun.
    \end{itemize}

\subsection{Quand utiliser \textit{de} vs \textit{du/de la/de l'/des}}
\textbf{Règle courte:} Utilise \textit{de} après une quantité, une négation, ou une expression figée; utilise \textit{du/de la/de l'/des} pour un article partitif ou indéfini selon le contexte.
\begin{center}
    \begin{tabular}{|l|l|l|}
        \hline
        \textbf{Forme} & \textbf{Usage} & \textbf{Exemple court} \\
        \hline
        de & quantité, négation & beaucoup de travail, pas de bruit \\
        \hline
        du/de la/de l' & partitif (quantité non comptée) & boire du café \\
        \hline
        des & pluriel indéfini & acheter des fruits \\
        \hline
        de (pas des) & adjectif pluriel avant le nom & de bons amis \\
        \hline
    \end{tabular}
\end{center}
\textbf{Exemple modèle:} \textit{Je prends du thé, mais je ne bois pas de café et j'achète des biscuits.} (English: I drink tea, but I don't drink coffee, and I buy cookies.) \\
\textbf{Piège fréquent:} Après une négation, on met généralement \textit{de} (\textit{pas de}), pas \textit{du/de la/des}.

\subsection{Quand utiliser \textit{qui} vs \textit{que} vs \textit{dont} vs \textit{où}}
\textbf{Règle courte:} Choisis le pronom selon la fonction dans la subordonnée: sujet, COD, complément avec \textit{de}, ou repère de lieu/temps.
\begin{center}
    \begin{tabular}{|l|l|l|}
        \hline
        \textbf{Forme} & \textbf{Usage} & \textbf{Exemple court} \\
        \hline
        qui & sujet du verbe suivant & La femme qui parle \\
        \hline
        que & COD du verbe suivant & Le film que j'ai vu \\
        \hline
        dont & verbe/adjectif + de & Le sujet dont on discute \\
        \hline
        où & lieu ou moment & Le jour où je suis parti \\
        \hline
    \end{tabular}
\end{center}
\textbf{Exemple modèle:} \textit{Le collègue qui arrive demain est celui que j'ai appelé et dont je t'ai parlé hier.} (English: The colleague arriving tomorrow is the one I called and told you about yesterday.) \\
\textbf{Piège fréquent:} Après \textit{de}, évite \textit{que}: on dit \textit{le livre dont je parle}, pas \textit{que je parle}.

\subsection{Quand utiliser \textit{ce qui} / \textit{ce que} / \textit{ce dont}}
These all mean something like ``what / which / the thing that'', but the choice depends on the role inside the following clause.

\textbf{\textit{ce qui}} \\
Use \textit{ce qui} when it is the subject of the next verb. \\
\textbf{Pattern:} \textit{ce qui + verb}
\begin{itemize}
    \item \textit{Ce qui m'intéresse, c'est la musique.} (English: What interests me is music.)
    \item \textit{Je ne sais pas ce qui se passe.} (English: I don't know what is happening.)
\end{itemize}
\textbf{Rule:} if the next verb still needs a subject, use \textit{ce qui}.

\textbf{\textit{ce que}} \\
Use \textit{ce que} when it is the direct object of the next verb. \\
\textbf{Pattern:} \textit{ce que + subject + verb}
\begin{itemize}
    \item \textit{Ce que j'aime, c'est la musique.} (English: What I like is music.)
    \item \textit{Je sais ce que tu veux.} (English: I know what you want.)
\end{itemize}
\textbf{Rule:} if the clause already has a subject and the verb needs an object, use \textit{ce que}.

\textbf{\textit{ce dont}} \\
Use \textit{ce dont} when the verb or expression requires \textit{de}. \\
\textbf{Pattern:} \textit{ce dont + subject + verb/expression with de}
\begin{itemize}
    \item \textit{Ce dont j'ai besoin, c'est de repos.} (English: What I need is rest.)
    \item \textit{Je comprends ce dont tu parles.} (English: I understand what you're talking about.)
\end{itemize}
\textbf{Common triggers:}
\begin{itemize}
    \item \textit{avoir besoin de}
    \item \textit{parler de}
    \item \textit{se souvenir de}
    \item \textit{rêver de}
\end{itemize}
\textbf{Rule:} if the verb/expression is built with \textit{de}, use \textit{ce dont}. \\

\textbf{Quick test}
\begin{itemize}
    \item \textit{ce qui} = missing subject
    \item \textit{ce que} = missing direct object
    \item \textit{ce dont} = missing \textit{de} + thing
\end{itemize}

\subsection{Quand utiliser \textit{lequel/laquelle/lesquels/lesquelles}}
\textbf{Règle courte:} Utilise la famille de \textit{lequel} surtout après une préposition pour reprendre une chose (ou un animal), avec accord en genre et en nombre.
\begin{center}
    \begin{tabular}{|l|l|l|}
        \hline
        \textbf{Forme} & \textbf{Usage} & \textbf{Exemple court} \\
        \hline
        lequel & masculin singulier & Le bureau sur lequel j'écris \\
        \hline
        laquelle & féminin singulier & La chaise sur laquelle je m'assois \\
        \hline
        lesquels & masculin pluriel & Les dossiers dans lesquels je classe \\
        \hline
        lesquelles & féminin pluriel & Les règles avec lesquelles j'étudie \\
        \hline
    \end{tabular}
\end{center}
\textbf{Exemple modèle:} \textit{Le dossier dans lequel j'ai mis les notes est sur la table sur laquelle tu travailles.} (English: The file in which I put the notes is on the table on which you are working.) \\
\textbf{Piège fréquent:} Pour une personne, on préfère souvent \textit{préposition + qui} (\textit{la collègue avec qui je travaille}) plutôt que \textit{avec laquelle}.

\subsection{Quand utiliser \textit{à qui}/\textit{auquel} et \textit{de qui}/\textit{duquel} (vs \textit{dont})}
\textbf{Règle courte:} Avec \textit{à} ou \textit{de}, le choix dépend de la nature du nom (personne/chose) et du type de préposition (simple ou composée).
\begin{center}
    \begin{tabular}{|l|l|l|}
        \hline
        \textbf{Forme} & \textbf{Usage} & \textbf{Exemple court} \\
        \hline
        à qui & personne avec \textit{à} & L'ami à qui j'écris \\
        \hline
        auquel... & chose avec \textit{à} & Le projet auquel je pense \\
        \hline
        de qui & personne avec \textit{de} & La personne de qui je parle \\
        \hline
        duquel... & chose avec \textit{de} après préposition & Le quartier près duquel j'habite \\
        \hline
        dont & remplace \textit{de} directement & Le livre dont j'ai besoin \\
        \hline
    \end{tabular}
\end{center}
\textbf{Exemple modèle:} \textit{La collègue à qui je parle travaille sur un sujet auquel je m'intéresse, et le quartier près duquel elle habite est calme.} (English: The colleague I am speaking to is working on a topic I am interested in, and the neighborhood near which she lives is quiet.) \\
\textbf{Piège fréquent:} Après une préposition composée avec \textit{de} (\textit{près de, à côté de}), on n'utilise pas \textit{dont}; on utilise \textit{duquel/de laquelle/de qui}.

\subsection{Quand utiliser \textit{à} vs \textit{en} vs \textit{dans} vs \textit{chez}}
\textbf{Règle courte:} Ces prépositions marquent des relations différentes de lieu, mouvement, contexte et appartenance sociale/personnelle.
\begin{center}
    \begin{tabular}{|l|l|l|}
        \hline
        \textbf{Forme} & \textbf{Usage} & \textbf{Exemple court} \\
        \hline
        à & ville, destination ponctuelle & aller à Paris \\
        \hline
        en & pays féminins/moyen de transport & vivre en France, venir en train \\
        \hline
        dans & intérieur, délai futur, inclusion & dans la maison, dans deux jours \\
        \hline
        chez & personne, professionnel, groupe & chez Paul, chez le médecin \\
        \hline
    \end{tabular}
\end{center}
\textbf{Exemple modèle:} \textit{Je vais à Paris, puis je reste en France et je dors chez une amie dans le centre-ville.} (English: I am going to Paris, then staying in France, and sleeping at a friend's place downtown.) \\
\textbf{Piège fréquent:} Ne confonds pas \textit{en} et \textit{au}: on dit \textit{en France}, mais \textit{au Canada}.

\subsection{Quand utiliser \textit{pour} vs \textit{par}}
\textbf{Règle courte:} \textit{pour} exprime surtout le but, la destination ou le bénéficiaire; \textit{par} exprime le moyen, la cause, le passage ou l'agent du passif.
\begin{center}
    \begin{tabular}{|l|l|l|}
        \hline
        \textbf{Forme} & \textbf{Usage} & \textbf{Exemple court} \\
        \hline
        pour & but, destinataire, échéance & étudier pour réussir \\
        \hline
        par & moyen, cause, trajet, agent passif & envoyé par mail, écrit par Hugo \\
        \hline
    \end{tabular}
\end{center}
\textbf{Exemple modèle:} \textit{Je travaille pour améliorer mon français et j'apprends par des podcasts recommandés par mon professeur.} (English: I work to improve my French, and I learn through podcasts recommended by my teacher.) \\
\textbf{Piège fréquent:} Dans une phrase passive, utilise \textit{par} pour l'agent (\textit{Ce texte est écrit par...}), pas \textit{pour}.

\subsection{Quand utiliser \textit{depuis} vs \textit{pendant} vs \textit{il y a} vs \textit{dans}}
\textbf{Règle courte:} Ces marqueurs temporels se distinguent par la relation au moment présent: durée en cours, durée terminée, passé ponctuel, futur.
\begin{center}
    \begin{tabular}{|l|l|l|}
        \hline
        \textbf{Forme} & \textbf{Usage} & \textbf{Exemple court} \\
        \hline
        depuis & durée commencée et encore vraie & J'étudie depuis 2023 \\
        \hline
        pendant & durée terminée (ou délimitée) & J'ai voyagé pendant un mois \\
        \hline
        il y a & point dans le passé (ago) & Je l'ai vu il y a deux jours \\
        \hline
        dans & moment futur à partir de maintenant & Je pars dans deux semaines \\
        \hline
    \end{tabular}
\end{center}
\textbf{Exemple modèle:} \textit{J'étudie le français depuis deux ans; j'ai fait une pause pendant un mois il y a trois semaines, et je reprendrai sérieusement dans deux jours.} (English: I've been studying French for two years; I took a one-month break three weeks ago, and I'll start seriously again in two days.) \\
\textbf{Piège fréquent:} N'utilise pas \textit{depuis} pour une action finie: avec un passé terminé, prends plutôt \textit{pendant}.

\subsection{Rule for de}

\begin{enumerate}
    \item If \textbf{de} comes after a \textbf{quantity expression}, use \textbf{de} alone.
    \item If \textbf{de} comes after a \textbf{negative expression}, indefinite/partitive articles usually become \textbf{de}.
    \item If \textbf{de} means \textbf{of/from/about/around/near} a specific noun, keep the article.
\end{enumerate}


\section{Volcab}
    \begin{itemize}
        \item passer pour (to be considered as, to pass for)
        \item se faire passer pour (to pretend to be)
        \item rendre visite à quelqu'un (for places, use ``visiter'')
        \begin{itemize}
            \item aller voir quelqu'un
            \item voir quelqu'un
        \end{itemize}
    \end{itemize}


\end{document}
